% ============================================================
\chapter{M\'ethodes Directes pour Syst\`emes Lin\'eaires}
\label{ch:directes}
% ============================================================

\section{Exemple motivant}
Un circuit \'electrique \`a $n$ n\oe uds produit un syst\`eme $Ax=b$ de $n$ \'equations. Pour $n=1000$, il faut une m\'ethode syst\'ematique et efficace.

\section{\'Elimination de Gauss et factorisation LU}

\begin{definition}[Factorisation LU]\index{factorisation LU}\index{LU}
Trouver $L$ triangulaire inf\'erieure (avec $\ell_{ii}=1$) et $U$ triangulaire sup\'erieure telles que $A=LU$. Alors $Ax=b$ se r\'esout par $Ly=b$ (descente) puis $Ux=y$ (remont\'ee).
\end{definition}

\begin{algorithm}[H]
\caption{Factorisation LU (sans pivot)}\label{alg:lu}
\KwEntree{Matrice $A\in\R^{n\times n}$}
\KwSortie{$L$, $U$ tels que $A=LU$}
\Pour{$k=1,\ldots,n-1$}{
  \Pour{$i=k+1,\ldots,n$}{
    $\ell_{ik} \gets a_{ik}/a_{kk}$\;
    \Pour{$j=k+1,\ldots,n$}{
      $a_{ij} \gets a_{ij}-\ell_{ik}\cdot a_{kj}$\;
    }
  }
}
$U \gets$ partie triangulaire sup\'erieure de $A$\;
\end{algorithm}

\begin{theoreme}[Existence de LU]
Si toutes les sous-matrices principales de $A$ sont inversibles, alors la factorisation LU existe et est unique.
\end{theoreme}

\begin{theoreme}[Co\^ut]
La factorisation LU co\^ute $\frac{2n^3}{3}+\BigO(n^2)$ op\'erations flottantes.
\end{theoreme}

\begin{proof}
L'\'etape $k$ effectue $(n-k)^2$ multiplications et soustractions. Total : $\sum_{k=1}^{n-1}2(n-k)^2=2\sum_{j=1}^{n-1}j^2=\frac{2n^3}{3}+\BigO(n^2)$.
\end{proof}

\section{Pivot partiel}\index{pivot partiel}

\begin{attention}
Sans pivotement, l'algorithme peut \^etre instable (division par un petit pivot). Le \textbf{pivot partiel} permute les lignes pour maximiser $|a_{kk}|$.
\end{attention}

On obtient $PA=LU$ o\`u $P$ est une matrice de permutation.

\begin{theoreme}[Stabilit\'e]
Avec pivot partiel, l'\'elimination de Gauss est \textbf{r\'etrograde stable} :
\[
(A+\Delta A)\hat{x}=b, \quad \|\Delta A\|\leq g(n)\,\eps_{\mathrm{mach}}\,\|A\|,
\]
o\`u $g(n)$ est le facteur de croissance (born\'e par $2^{n-1}$ en th\'eorie, rarement d\'epass\'e en pratique).
\end{theoreme}

\section{Factorisation de Cholesky}\index{Cholesky}

\begin{theoreme}[Factorisation de Cholesky]
Si $A$ est \textbf{sym\'etrique d\'efinie positive} (SDP), alors il existe une unique matrice triangulaire inf\'erieure $L$ \`a coefficients diagonaux positifs telle que $A=LL^\top$.
\end{theoreme}

\begin{proof}
Par r\'ecurrence sur $n$. \'Ecrivons $A=\begin{pmatrix}a_{11}&\mathbf{v}^\top\\\mathbf{v}&A'\end{pmatrix}$. Comme $A$ SDP, $a_{11}>0$. Posons $\ell_{11}=\sqrt{a_{11}}$, $\mathbf{l}=\mathbf{v}/\ell_{11}$. Alors $A'- \mathbf{l}\mathbf{l}^\top$ est SDP de taille $(n-1)\times(n-1)$ et on applique l'hypoth\`ese de r\'ecurrence.
\end{proof}

\begin{algorithm}[H]
\caption{Factorisation de Cholesky}\label{alg:cholesky}
\KwEntree{$A$ SDP, $n\times n$}
\KwSortie{$L$ telle que $A=LL^\top$}
\Pour{$j=1,\ldots,n$}{
  $\ell_{jj}\gets\sqrt{a_{jj}-\sum_{k=1}^{j-1}\ell_{jk}^2}$\;
  \Pour{$i=j+1,\ldots,n$}{
    $\ell_{ij}\gets\frac{1}{\ell_{jj}}\left(a_{ij}-\sum_{k=1}^{j-1}\ell_{ik}\ell_{jk}\right)$\;
  }
}
\end{algorithm}

\begin{theoreme}[Co\^ut de Cholesky]
$\frac{n^3}{3}+\BigO(n^2)$ : deux fois moins que LU.
\end{theoreme}

\section{Exemple num\'erique}

$A=\begin{pmatrix}4&2\\2&5\end{pmatrix}$, $b=\begin{pmatrix}8\\9\end{pmatrix}$. Cholesky : $L=\begin{pmatrix}2&0\\1&2\end{pmatrix}$.

Descente $Ly=b$ : $y_1=4$, $y_2=(9-1\cdot4)/2=2{,}5$. Remont\'ee $L^\top x=y$ : $x_2=2{,}5/2=1{,}25$, $x_1=(4-1\cdot 1{,}25)/2=1{,}375$.

\section{Impl\'ementation}

\begin{codeblock}[title=\textbf{Python}]
\begin{minted}{python}
import numpy as np
from scipy.linalg import lu, cholesky, solve_triangular

# LU
A = np.array([[2., 1, 1], [4, 3, 3], [8, 7, 9]])
b = np.array([4., 10, 24])
P, L, U = lu(A)
y = solve_triangular(L, P @ b, lower=True)
x = solve_triangular(U, y)
print(f"LU: x = {x}")

# Cholesky
A_spd = np.array([[4., 2], [2, 5]])
b_spd = np.array([8., 9])
L_chol = cholesky(A_spd, lower=True)
y = solve_triangular(L_chol, b_spd, lower=True)
x = solve_triangular(L_chol.T, y)
print(f"Cholesky: x = {x}")
print(f"Cout LU (n=100): ~{2*100**3//3:.0f} flops")
\end{minted}
\end{codeblock}

\begin{codeblock}[title=\textbf{Julia}]
\begin{minted}{julia}
using LinearAlgebra
A = [2.0 1 1; 4 3 3; 8 7 9]
b = [4.0, 10, 24]
F = lu(A)
x = F \ b
println("LU: x = $x")

A_spd = [4.0 2; 2 5]
b_spd = [8.0, 9]
C = cholesky(A_spd)
x = C \ b_spd
println("Cholesky: x = $x")
\end{minted}
\end{codeblock}

\section{Exercices}

\begin{exercice}[$\star$]
Factoriser $A=\begin{pmatrix}1&2&3\\2&8&14\\3&14&34\end{pmatrix}$ en $LU$ \`a la main. V\'erifier.
\end{exercice}

\begin{exercice}[$\star\star$]
Montrer que si $A$ est SDP, tous ses pivots dans l'\'elimination de Gauss sont positifs (sans pivotement n\'ecessaire).
\end{exercice}

\begin{exercice}[$\star\star$]
Comparer les temps de calcul de LU vs.\ Cholesky pour $n=100,500,1000,2000$ sur des matrices SDP al\'eatoires. V\'erifier le rapport $\sim 2$.
\end{exercice}

\begin{exercice}[$\star\star\star$ -- Projet]
Impl\'ementer la factorisation LU bande pour des matrices tridiagonales (algorithme de Thomas). Appliquer \`a la discr\'etisation de $-u''=f$.
\end{exercice}

\begin{formules}
\begin{align*}
A &= LU \quad (\text{co\^ut } \tfrac{2n^3}{3}) \\
A &= LL^\top \quad (\text{Cholesky, co\^ut } \tfrac{n^3}{3}) \\
PA &= LU \quad (\text{pivot partiel}) \\
\cond(A) &= \|A\|\cdot\|A^{-1}\|, \quad \frac{\|\delta x\|}{\|x\|}\leq\cond(A)\frac{\|\delta b\|}{\|b\|}
\end{align*}
\end{formules}
